\section{Compilation into finite numerical tables}\label{sec:effective}
A covering theorem supplies a finite set, but does not in general supply
an algorithm for choosing that set. Moreover, a table with finitely many
entries may still contain noncomputable real numbers. We separate these
issues. The construction below uses finite grids of feasible histories,
ordinary rational comparisons, and approximations to a finite moment
table. Neither rank detection nor a decision whether two exponents
coincide is among its operations.

\subsection{A compilation theorem for compact state systems}
Let $T\ge1$ be fixed, let $U=[\eta,1-\eta]^J$, with rational
$0<\eta<1/2$, and let $\mathcal X$ be finite. For $0\le n\le T$
let $S_n\subset[-B,B]^{d_n}$ be compact, with $S_0$ a singleton.
There are transition maps
\[
 T_{n,x}:S_n\times U\longrightarrow S_{n+1},\qquad
 S_{n+1}=\bigcup_{x\in\mathcal X}T_{n,x}(S_n,U).
\]
Norms in this section are maximum norms. Suppose
\begin{equation}\label{eq:compiler-lipschitz}
 \|T_{n,x}(s,u)-T_{n,x}(s',u')\|_\infty
 \le L_n\|s-s'\|_\infty+G_n\|u-u'\|_\infty.
\end{equation}
Each stage has finitely many queries, with a prediction map
$Q_n:S_n\to[0,1]^{b_n}$ of Lipschitz constant $K_n$.
The bounds and all dimensions are fixed, independent of the state budget.
Write $H_n(w,u_1,\ldots,u_n)$ for the exact state of the report word
$w\in\mathcal X^n$. Iteration of \eqref{eq:compiler-lipschitz} shows that
its command dependence is $A_n$-Lipschitz, where
\[
 A_0=0,\qquad A_{n+1}=L_nA_n+G_n.
\]
For $M\ge1$, define the unrestricted-centre covering radius
\[
 e_n(M)=\inf_{z_1,\ldots,z_M\in\mathbb R^{d_n}}
                  \sup_{s\in S_n}\min_i\|s-z_i\|_\infty.
\]
No spectral representation or dimension formula for $e_n(M)$ is assumed.

\begin{definition}[Finite evaluation data]\label{def:finite-data}
Given a rational grid history and a positive rational tolerance $\tau$,
an evaluation procedure returns a dyadic vector within $\tau$ of its
exact state. It also returns dyadic query predictions within $\tau$ of
the exact predictions at any chosen representative history. This is an
offline procedure. In its effective version it is a specified algorithm
with access to the numerical input oracles of the experiment. In a
noneffective version the finitely many approximations are supplied data.
Neither version gives an oracle to the running transducer.
\end{definition}

Let $U_h\subset U$ be a finite rational $h$-net in the maximum norm,
with $|U_h|\le C_U(1+h^{-1})^J$. At runtime an input includes an index
of any $\bar u\in U_h$ satisfying $\|u-\bar u\|_\infty\le h$.
The guarantee below holds for every such choice. In particular it does
not require deciding equality of an arbitrary real input to a rounding
boundary. An input approximation, followed by rounding that rational
approximation and clipping to $U$, supplies this interface with a fixed
constant adjustment of $h$. For comparison with a memory lower bound,
input coding is a fixed memoryless Borel function of the current command,
or uses fresh independent rounding coins. It is not an additional
history-dependent channel. The upper estimate itself holds for every
name meeting the error bound.

\begin{lemma}[Greedy covering from approximate samples]\label{lem:finite-greedy}
Let $S$ be a compact subset of a normed space. Suppose a finite indexed
list of points $s_j\in S$ is an $a$-net of $S$, and
$\|\widetilde s_j-s_j\|\le\tau$. Farthest-first selection of at most
$M$ points from the rational list $(\widetilde s_j)$, with fixed tie
breaking, selects indices $I$ for which
\[
 \sup_{s\in S}\min_{i\in I}\|s-s_i\|
                 \le 2e_S(M)+a+4\tau.
\]
If there are at most $M$ indices, select all of them. Distinct indices
may have identical approximate vectors.
\end{lemma}
\begin{proof}
Let $r$ be the final covering radius of the approximate list. When the
list has more than $M$ indices, the selected $M$ points and a farthest
point have mutual distances at least $r$. Every cover by $M$ balls of
radius $b$ must put two of these $M+1$ points in the same ball, and
therefore $r\le2b$. This also handles $r=0$. A cover of $S$ with
radius $e_S(M)+\epsilon$ covers the approximate list with radius
$e_S(M)+\epsilon+\tau$, allowing its centres to lie outside $S$.
Thus $r\le2e_S(M)+2\tau$ after $\epsilon\downarrow0$.
For a given $s\in S$, choose $s_j$ within $a$, and then a selected
approximate point within $r$ of $\widetilde s_j$. The triangle
inequality gives $a+2\tau+r$. The case of a list of size at most $M$
is immediate. This is the classical farthest-first argument
\cite{Gonzalez}; the approximation allowance is recorded explicitly.
\end{proof}

\begin{theorem}[Effective compilation without loss of covering order]
\label{thm:compact-compiler}
Under \eqref{eq:compiler-lipschitz} and Definition~\ref{def:finite-data},
there is a finite procedure producing stage-dependent integer transition
tables and dyadic output tables, with at most $M$ state indices at each
stage, such that, on every input-report history,
\begin{equation}\label{eq:compiler-bound}
 \|\widehat Q_n-Q_n(s_n)\|_\infty
 \le C_T\left(\max_{1\le j\le n}e_j(M)+h+\tau\right).
\end{equation}
The running procedure reads only its previous index, the current finite
input code and report, and, at a checkpoint, the new query label. It
stores no numerical state vector, representative history, or oracle.
The constant depends only on the displayed bounds and the fixed horizon.
The synthesis procedure does not need $e_n(M)$ or a covering certificate
as input.
\end{theorem}
\begin{proof}
Enumerate all pairs of a report word in $\mathcal X^n$ and a command
word in $U_h^n$. They are feasible histories by the hypotheses on the
state system. Their exact images form an $A_nh$-net of $S_n$.
Evaluate every image with error at most $\tau$ and apply
Lemma~\ref{lem:finite-greedy}. Associate each selected index with the
\emph{exact} state $c_{n,i}$ of its selected history, for analysis only.
The stored approximations $\widetilde c_{n,i}$ have error at most
$\tau$. The exact representatives cover $S_n$ with radius
\begin{equation}\label{eq:compiled-cover}
 R_n\le2e_n(M)+A_nh+4\tau.
\end{equation}

For each representative history at stage $n$, each $\bar u\in U_h$
and each $x$, append $(\bar u,x)$ to that history and evaluate the
resulting exact state
$y=T_{n,x}(c_{n,i},\bar u)$ with error at most $\tau$.
Choose the next index by minimizing the maximum-norm distance between
this approximation and the stored approximations at stage $n+1$.
All comparisons are between rational numbers. If $j$ is selected, the
triangle inequality, applied to a nearest exact centre, gives
\[
 \|y-c_{n+1,j}\|_\infty\le R_{n+1}+4\tau.
\]
This specifies one integer for every entry of a finite transition table.
The offline representative history is not a runtime variable. In
particular no update is applied to an approximate vector outside the
reachable set: $y$ is defined and evaluated through an actual history.

Let $E_n=\|s_n-c_{n,i_n}\|_\infty$, where $s_n$ denotes the actual
uncompressed state, again only in the analysis. For any admitted finite
input code,
\begin{equation}\label{eq:digital-recurrence}
 E_0=0,\qquad
 E_{n+1}\le L_nE_n+G_nh+2e_{n+1}(M)+A_{n+1}h+8\tau.
\end{equation}
Induction over the fixed horizon bounds $E_n$ by a constant times the
right-hand quantity in \eqref{eq:compiler-bound}. Store a dyadic
approximation to $Q_n(c_{n,i})$ with error at most $\tau$ and clip it
to $[0,1]$. Clipping cannot increase its error to a probability.
The query error is at most $K_nE_n+\tau$, proving the theorem.
There are finitely many evaluations and comparisons, so the procedure
terminates whenever the evaluation procedure does. Nonuniqueness of
representatives or ties does not require an equality test on exact data.
\end{proof}

\begin{proposition}[Size of the finite program]\label{prop:compiler-size}
Put $V=|U_h|$ and $X=|\mathcal X|$. The transition table has at most
$TMVX$ entries, each taking $\lceil\log_2(M+1)\rceil$ bits. If outputs
are written to $b$ binary places, the output tables take at most
$M\sum_nb_n(b+1)$ bits, in addition to fixed metadata. Candidate
enumeration uses at most $\sum_{n=0}^T(XV)^n$ histories. With cached
candidate approximations, greedy selection uses
$O(M\sum_{n=0}^T(XV)^n)$ fixed-dimensional distance operations, and
transition selection uses $O(TM^2VX)$ such operations.
\end{proposition}
\begin{proof}
The transition arguments are the stage, current representative, finite
command code and report. Counting their combinations proves the first
assertion. A query table has one $b$-place probability per representative
and query. At stage $n$ there are $(XV)^n$ grid histories. Maintain for
each candidate its smallest distance to a chosen centre; one new centre
updates this array with one pass, repeated at most $M$ times. For each
transition, comparison against at most $M$ next centres gives the last
count. Exact integer comparison, maximum and subtraction suffice.
These counts exclude the running time of an arbitrary numerical oracle.
\end{proof}

With a supplied clock the only persistent field has at most $M$ values.
An autonomous machine can retain the pair $(n,i_n)$, using at most
$1+(T+1)M$ values. The current input code, table address and returned
output are transient data. They require $O(\log M+\log V+b+\log(T+1))$
bits for fixed dimensions. The finite read-only program is a separate,
explicitly counted resource; it is not claimed to fit into $\log M$ bits.

\subsection{Finite data for positive monomial experiments}
Return to the experiment of Section~\ref{sec:collision-geometry}, with
$I=[0,1]$, and put $T=N-1$. Use the full formal raw moment vector
$v_n=(\E[t^{\alpha\cdot a}\mid h])_{|\alpha|=N-n}$, including its
constant entry. Coincident labels remain present. Its dimension is
bounded by a constant depending only on $r,N$.

Choose a fixed bound
\[
 B_f\ge1+\sup_{g,x}\sum_{i=0}^{r-1}|f_i(g,x)|,
 \qquad f_{g,x}(t)=\sum_i f_i(g,x)t^{a_i}.
\]
This uses the coefficient realization, not exponent separation.
For two reachable raw vectors, the update numerator and denominator
have differences at most $B_f\|v-v'\|_\infty$. Their denominators
are at least $\kappa$, and a numerator divided by its own denominator
belongs to $[0,1]$. Subtracting the two quotients gives a state
Lipschitz bound $2B_f/\kappa$. For fixed posterior, the difference of
report factors at two commands is at most $J\|g-g'\|_\infty$
pointwise, so the command bound is $2J/\kappa$. Consequently
\eqref{eq:compiler-lipschitz} holds uniformly on the calibration set.
Physical query predictions are fixed linear combinations of raw moments,
with uniformly bounded coefficients. For effective evaluation we fix
a sufficiently small positive rational $\delta$ in the canonical
probe basis $1/2$, $1/2+\delta t^{a_i}$. Its formal query coefficients
are then rational. Alternatively certified approximations to a chosen
fixed query coefficient table are included among the numerical inputs.
This verifies every metric
hypothesis of Theorem~\ref{thm:compact-compiler}.

\begin{lemma}[Certified finite evaluation from prior moments]
\label{lem:finite-moments}
For a fixed horizon, dyadic approximations to the coefficients of the
cell matrix and to
\begin{equation}\label{eq:finite-prior-table}
 m_\gamma(a)=\int_0^1t^{\gamma\cdot a}\,d\mu(t),\qquad
 \gamma\in\mathbb N_0^r,\quad |\gamma|=N,
\end{equation}
suffice to evaluate every rational-command history state and every fixed
finite product-query prediction with a prescribed error $\tau>0$.
Required input accuracy is $c_N\min\{\tau,1\}$ for a positive
constant determined by supplied coefficient and evidence bounds. The
number of input moment entries is independent of the history grid and
of $M$. No equality decision on the exponents is required.
\end{lemma}
\begin{proof}
Expand an $n$-report product in formal homogeneous variables as
$P(X)=\sum_{|\beta|=n}p_\beta X^\beta$. Its coefficient sum satisfies
$\sum_\beta|p_\beta|\le B_f^n$. A coordinate of the remaining raw
state is
\begin{equation}\label{eq:finite-evaluation}
 v_\alpha=
 \frac{\sum_{|\beta|=n}p_\beta m_{\alpha+\beta}}
      {\sum_{|\beta|=n}p_\beta m_{\beta+(N-n)e_0}},
 \qquad |\alpha|=N-n.
\end{equation}
The denominator is the true prefix evidence and is at least $\kappa^n$.
The formula remains exact when distinct formal labels have equal
exponents. There is no cancellation-based inference of rank.

A perturbation of size $\delta$ in the input coefficients changes each
factor coefficient vector by at most a fixed multiple of $\delta$.
The telescoping identity for products bounds the coefficient-sum error
in $P$ by $n(B_f+1)^{n-1}C\delta$, for $\delta$ sufficiently small.
Each prior moment lies in $[0,1]$, and there are finitely many formal
labels. Thus numerator and denominator errors are at most $D_N\delta$,
where $D_N$ is obtained from the finite bounds just displayed.
Take $D_N\delta\le\kappa^N/2$. The approximate denominator is then
positive. If a true numerator is $A$, its denominator is $Z$, and
$0\le A\le Z$, direct subtraction gives
\[
 \left|\frac{\widetilde A}{\widetilde Z}-\frac AZ\right|
 \le\frac{|\widetilde A-A|+|\widetilde Z-Z|}{\widetilde Z}
 \le4D_N\kappa^{-N}\delta.
\]
Choose $\delta$ to make this at most $\tau/2$ and round the result
to a dyadic within $\tau/2$. Fixed product-query predictions are
linear combinations of these raw entries, or the same product ratios;
increase the accuracy constant by their fixed coefficient bounds.
This gives the claimed finite procedure.
\end{proof}

\begin{theorem}[Finite-precision realization of the intrinsic law]
\label{thm:digital-intrinsic}
Fix the experiment, the full-support prior, the horizon and compact
calibration set of Section~\ref{sec:collision-geometry}. Suppose the
command cube is rationally specified, and positive rational lower bounds
for the report likelihoods and positive exponents, together with upper
coefficient bounds, are supplied. For each known calibration $a$ and
integer $M\ge1$, finite approximations to the data in
Lemma~\ref{lem:finite-moments} produce a deterministic finite-table
transducer with at most $M$ persistent labels per stage and
\begin{equation}\label{eq:digital-intrinsic}
 \sup_{n,\mathfrak h}\mathcal R(\widehat Q_n,Q_n(\mathfrak h))
 \le C\{\Xi_N(M,a)+h^2+\tau^2\}.
\end{equation}
Here $h$ is the command mesh and $\tau$ the certified error of the
offline state and output evaluations. Constants are uniform in $a,M$.
For meshes and evaluation errors at most $c/(M+1)$, this transducer
has regret at most $C'\Xi_N(M,a)$, while every such transducer with the memoryless input interface has
maximum unconditional checkpoint regret at least $c'\Xi_N(M,a)$
under the same exploration law as in the intrinsic theorem.

All numerical input and output precisions can be chosen
$\log_2(M+1)+O(1)$ binary places. The transition and output program has
$O(M^{J+1}\log(M+1))$ bits for the fixed experiment. Its synthesis
uses polynomially many fixed-dimensional rational operations in $M$,
with exponent allowed to depend on $J,N$. If the coefficients and
moments have computable approximation procedures, the synthesis is a
single uniform terminating algorithm using those procedures offline.
No computability assumption is needed for the preceding existential
intrinsic law for arbitrary full-support priors.
\end{theorem}
\begin{proof}
Theorem~\ref{thm:intrinsic-checkpoint}, its reachable-centre construction,
and the fixed physical/raw metric comparison imply
\[
 e_n(M)\le C_n\Psi_{n,N-n}(M,a)^{1/2}.
\]
The raw update calculation above and Lemma~\ref{lem:finite-moments}
verify all assumptions of Theorem~\ref{thm:compact-compiler}. Squaring
\eqref{eq:compiler-bound}, using
$(u+v+w)^2\le3(u^2+v^2+w^2)$, and averaging over the finite query
menu prove \eqref{eq:digital-intrinsic}. For $N\ge2$, the $\ell=1$
term at every nontrivial checkpoint gives
\[
 \Xi_N(M,a)\ge M^{-2}.
\]
Hence command and evaluation errors of order $(M+1)^{-1}$ are absorbed
without determining any individual pivot or collision scale.

A dyadic grid with this mesh has $O(M^J)$ points. Coefficient and prior
moment errors of size $c_N/(M+1)$ suffice by
Lemma~\ref{lem:finite-moments}; the required binary precision is the
stated logarithmic quantity. Proposition~\ref{prop:compiler-size} gives
the storage and arithmetic counts. Dyadic input coefficients may be
converted to rational product ratios of $O_N(\log(M+1))$ bits, and all
stored approximations have that precision. Integer comparisons are
finite, including exact ties.

Finally the finite-table filter, even when its command input is a
quantized approximation, is an admissible filter with at most $M$
labels in the original experiment when input coding is memoryless.
The exact-input filter can first apply that fixed coding rule. It has less, not more, information
than one reading the exact command. The lower bound in
Theorem~\ref{thm:intrinsic-streaming} therefore applies unchanged.
This is a lower bound for one filter over the common exploration law,
not a combination of independently redesigned checkpoint filters.
\end{proof}

\begin{corollary}[Calibration accuracy independent of collision gaps]
\label{cor:calibration-precision}
Suppose $a_i\ge a_*>0$ for $i>0$, the calibration entries are
computable reals, and the prior has a procedure that approximates
$\int t^b\,d\mu$ for positive rational $b$ to requested accuracy.
The preceding compiler needs only $\log_2(M+1)+O(1)$ binary places
of the calibration and of the resulting moments. The bound does not
involve the smallest nonzero additive exponent gap and is valid at
exact collisions.
\end{corollary}
\begin{proof}
For $b,b'\ge a_*/2$, the mean value theorem and
$\sup_{0<t\le1}t^{a_*/2}|\log t|=2/(e a_*)$ give
\[
 \left|\int(t^b-t^{b'})\,d\mu\right|
                   \le\frac{2}{e a_*}|b-b'|.
\]
Every nonconstant exponent in \eqref{eq:finite-prior-table} is at least
$a_*$. Approximate its constituent calibration entries so that its
error is $O((M+1)^{-1})$ and below $a_*/2$. Evaluate the corresponding
positive rational exponent to the same accuracy. The constant entry
is exactly one. There are finitely many labels, so the constants can
be chosen uniformly. This argument neither compares two sums nor asks
whether their difference vanishes.
\end{proof}

\begin{corollary}[Digital realization of high-contact phases]
\label{cor:digital-contact}
For the two-parameter family in Corollary~\ref{cor:two-parameter}, the
same three-term regret law is attained by finite numerical transition
and output tables under the finite-data assumptions above. Along
$u=\theta$, $v=\theta+\theta^k$, the required numerical precision is
$\log_2(M+1)+O(1)$, with no extra term proportional to
$\log(1/|v-u|)$. In particular the compiler remains valid before that
smallest separation can be numerically resolved.
\end{corollary}
\begin{proof}
Apply Theorem~\ref{thm:digital-intrinsic} and
Corollary~\ref{cor:calibration-precision}. The three-term expression and
its contact orders are those of Corollary~\ref{cor:two-parameter}.
The compiler never identifies the three collision lines or computes
$k$; its accuracy is set by $M$, and its performance follows from the
complete attainable cover. The physical experiment and its full-history
loss baseline are unchanged.
\end{proof}

The construction is not a new interpolation theorem or a new
farthest-first approximation algorithm. It connects a collision-uniform
attainability theorem to a finite numerical state machine with a counted
program. The distinction matters: a purely spectral identity supplies
neither the attainable cover nor the positive-likelihood update bound.
The horizon is fixed throughout. The program can be much larger than
the persistent state, and the running time of a general moment oracle
has not been bounded. For a noncomputable prior, finite advice exists
at each precision but a uniform algorithm for producing that advice is
not asserted. These distinctions leave the arbitrary-prior mathematical
classification intact.
